\chapter*{参考文献}
\phantomsection
\addcontentsline{toc}{chapter}{参考文献}

\begin{enumerate}
\def\labelenumi{(\Roman{enumi})}
\item
  L. EULER, Opera Omnia: De formulas integralibus duplicatis (1), v. XVII, Leipzig--Berlin (Teubner), 1915, pp.~289--315.
\item
  G. LEJEUNE--DIRICHLET, Sur la convergence des séries trigonométriques qui servent à représenter une fonction arbitraire entre des limites données, J. de Crelle, 4 (1829), pp.~157--169 (= Werke, v. I, pp.~118--132, Berlin (G. Reimer), 1889).
\item
  B. RIEMANN, Gesammelte Mathematische Werke, 2nd. edn., Leipzig (Teubner), 1892.
\item
  G. CANTOR, Gesammelte Abhandlungen, Berlin (Springer), 1932.
\item
  G. PEANO, Applicazioni geometriche del calcolo infinitesimale, Turin, 1887.
\item
  C. JORDAN, Cours d'Analyse, v. I, 2nd. edn., Paris (Gauthier-Villars), 1893.
\item
  C. ARZELÀ: a) Sulla integrabilità di una serie di funzioni, Rendic. Acc. dei Lincei, (4), 1 (1885), pp.~321--326; b) Sulla integrazione per serie, ibid., pp.~532--537 and 566--569.
\item
  T. STIELTJES, Recherches sur les fractions continues, Ann. Fac. Sci. de Toulouse, 8 (1894), J. 1 to J. 122.
\item
  E. BOREL, Leçons sur la théorie des fonctions, Paris (Gauthier-Villars), 1898.
\item
  H. LEBESGUE: a) Intégrale, longueur, aire, Annali di Mat., (3), 7 (1902), pp.~231--359; b) Leçons sur l'Intégration et la recherche des fonctions primitives, Paris (Gauthiers--Villars), 1904; c) Sur les séries trigonométriques, Ann. Éc. Norm. Sup., (3), 20 (1903), pp.~453--485; d) Sur l'intégration des fonctions discontinues, Ann. Éc. Norm. Sup., (3), 27 (1910), pp.~361--450.
\item
  W. H. YOUNG: a) On upper and lower integration, Proc. Lond. Math. Soc., (2), 2 (1905), pp.~52--66; b) A new method in the theory of integration, Proc. Lond. Math. Soc., (2), 9 (1911), pp.~15--50.
\item
  G. VITALI: a) Una proprieta delle funzioni misurabili, R. Ist. Lombardo, Rendiconti, (2), 38 (1905), pp.~599--603; b) Sui gruppi di punti e sulle funzioni di variabili reali, Rendic. Acc. Sci. di Torino, 43 (1908), pp.~229--236.
\item
  E. FISCHER, Sur la convergence en moyenne, C. R. Acad. Sci., 144 (1907), pp.~1022--1024.
\item
  G. FUBINI, Sugli integrali multipli, Rendic. Acc. dei Lincei, (5), 16 (1907), pp.~608--614.
\item
  F. RIESZ: a) Sur les systèmes orthogonaux de fonctions, C. R. Acad. Sci., 144 (1907), pp.~615--619; b) Untersuchungen über Systeme integrierbarer Funktionen, Math. Ann., 69 (1910), pp.~449--497; c) Sur certains systèmes singuliers d'équations intégrales, Ann. Éc. Norm. Sup., (3), 28 (1911), pp.~33--62; d) Sur quelques notions fondamentales dans la théorie générale des opérations linéaires, Ann. of Math., (2), 41 (1940), pp.~174--206.
\item
  D. EGOROFF, Sur les suites de fonctions mesurables, C. R. Acad. Sci., 152 (1911), p.~244.
\item
  J. RADON, Theorie und Anwendungen der absolut additiven Mengenfunktionen, Sitzungsber. der math. naturwiss. Klasse der Akad. der Wiss. (Wien), 122, Abt. IIa (1913), pp.~1295--1438.
\item
  C. CARATHÉODORY, Vorlesungen über reelle Funktionen, Leipzig--Berlin (Teubner), 1918.
\item
  P. J. DANIELL: a) A general form of integral, Ann. of Math., (2), 19 (1918), pp.~279--294; b) Integrals in an infinite number of dimensions, Ann. of Math., (2), 20 (1919), pp.~281--288.
\item
  O. NIKODYM, Sur une généralisation des intégrales de M. J. Radon, Fund. Math., 15 (1930), pp.~131--179.
\item
  S. BOCHNER, Integration von Funktionen deren Werte die Elemente eines Vektorraumes sind, Fund. Math., 20 (1933), pp.~262--276.
\item
  J. von NEUMANN, On rings of operators. III, Ann. of Math., (2), 41 (1940), pp.~94--161.
\item
  A. KOLMOGOROFF, Grundbegriffe der Wahrscheinlichkeitsrechnung, Berlin (Springer), 1933.
\item
  S. SAKS, Theory of the integral, 2nd. edn., New York (Stechert), 1937.
\item
  A. ZYGMUND, Trigonometric series, Warsaw, 1935 (2nd. edn., Cambridge University Press, 1959).
\item
  L. CESARI, Surface area, Princeton, 1954.
\item
  L. H. LOOMIS, An introduction to abstract harmonic analysis, London--New York--Toronto (van Nostrand), 1953.
\end{enumerate}

